\mxnewpage\section*{Correction EV2-T2}\label{EV2-T2-CORR}
\textbf{QCM :} 1 D, 2 D, 3 A, 4 C, 5 C, 6 D, 7 D, 8 D, 9 B, 10 A, 11 D, 12 B. Chaque réponse correcte : 0{,}5 point ; total 6.
\paragraph{Vérifications des automatismes.} Priorités ; somme de fractions ; soustraction des exposants ; coefficient multiplicateur 1{,}2 ; division par 0{,}8 ; produit 1{,}25 fois 0{,}8 ; résolution affine ; accroissement vertical divisé par l’horizontal ; moyenne ; angle divisé par 360 ; complément à 1 ; substitution entre parenthèses.
\subsection*{Exercice 1 -- Deux tirages}
\begin{enumerate}
\item $(\dfrac{4}{6})(\dfrac{3}{5})=\dfrac{2}{5}$. \textit{Barème : 1 point. Résultat et justification élémentaire attendus ; demi-point possible pour démarche correcte inachevée.}
\item $(\dfrac{4}{6})(\dfrac{2}{5})+(\dfrac{2}{6})(\dfrac{4}{5})=\dfrac{8}{15}$. \textit{Barème : 2 points. Méthode ou justification : 1 ; calcul et conclusion : 1.}
\item Avec remise : $2(\dfrac{2}{3})(\dfrac{1}{3})=\dfrac{4}{9}$. Les modèles diffèrent ; $\dfrac{8}{15}>\dfrac{4}{9}$. \textit{Barème : 2 points. Méthode ou justification : 1 ; calcul et conclusion : 1.}
\end{enumerate}
\subsection*{Exercice 2 -- Une boîte}
\begin{enumerate}
\item $V\prime=4(5-x)(5-3x)$. \textit{Barème : 1 point. Résultat et justification élémentaire attendus ; demi-point possible pour démarche correcte inachevée.}
\item Croît jusqu’à $\dfrac{5}{3}$ puis décroît jusqu’à 5. \textit{Barème : 2 points. Méthode ou justification : 1 ; calcul et conclusion : 1.}
\item $V(\dfrac{5}{3})=\dfrac{2000}{27}$ ; valeurs nulles aux bornes. \textit{Barème : 2 points. Méthode ou justification : 1 ; calcul et conclusion : 1.}
\end{enumerate}
\subsection*{Exercice 3 -- Reconnaissance et intersection}
\begin{enumerate}
\item $(x+2)^2+(y-1)^2=9$ : centre (-2;1), rayon 3. \textit{Barème : 1 point. Résultat et justification élémentaire attendus ; demi-point possible pour démarche correcte inachevée.}
\item Points $(-2;4)$ et $(-2;-2)$. \textit{Barème : 1 point. Résultat et justification élémentaire attendus ; demi-point possible pour démarche correcte inachevée.}
\item Rayon vertical au point supérieur, tangente horizontale $y=4$. Normale (0;1). \textit{Barème : 2 points. Méthode ou justification : 1 ; calcul et conclusion : 1.}
\end{enumerate}
\paragraph{Étalonnage.} Durée cible : 20–25 min pour les automatismes, 80–85 min pour les exercices, 10–15 min de relecture. Ces durées sont des estimations éditoriales à mesurer en classe. Accepter toute méthode correcte ; ne pas sanctionner deux fois une erreur numérique propagée avec cohérence. Le barème proposé ne vaut pas décision d’harmonisation de l’équipe.
